
\documentclass[11pt,a4paper]{article}

\usepackage[margin=2.45cm]{geometry}
\usepackage{amsmath,amssymb,amsthm,mathtools}
\usepackage{booktabs}
\usepackage{array}
\usepackage{microtype}
\usepackage[hidelinks]{hyperref}
\usepackage{xcolor}
\usepackage{enumitem}

\newtheorem{theorem}{Theorem}[section]
\newtheorem{proposition}[theorem]{Proposition}
\newtheorem{lemma}[theorem]{Lemma}
\newtheorem{corollary}[theorem]{Corollary}
\newtheorem{conjecture}[theorem]{Conjecture}
\newtheorem{definition}[theorem]{Definition}
\newtheorem{remark}[theorem]{Remark}
\newtheorem{problem}[theorem]{Open Problem}

\newcommand{\E}{\mathbb E}
\newcommand{\PP}{\mathbb P}
\newcommand{\Law}{\mathrm{Law}}
\newcommand{\Var}{\mathrm{Var}}
\newcommand{\Cov}{\mathrm{Cov}}
\newcommand{\TV}{\mathrm{TV}}
\newcommand{\W}{W}
\newcommand{\cB}{\mathcal B}
\newcommand{\cC}{\mathcal C}
\newcommand{\cD}{\mathcal D}
\newcommand{\cE}{\mathcal E}
\newcommand{\cK}{\mathcal K}
\newcommand{\cL}{\mathcal L}
\newcommand{\bits}{\{0,1\}}
\newcommand{\supp}{\operatorname{supp}}
\newcommand{\one}{\mathbf 1}
\newcommand{\cumul}{\kappa}
\newcommand{\ip}{\imath}
\newcommand{\eps}{\varepsilon}
\newcommand{\R}{\mathbb R}
\newcommand{\Z}{\mathbb Z}

\title{\textbf{Who Pays the Bill?}\\[1mm]
\large Checkpoint v6: Sharp Second Order, Higher-Order Causal Buffers,\\
Parity Ladders, and the Zero-Repair Viability Problem}
\author{}
\date{21 August 2026}

\begin{document}
\maketitle

\begin{abstract}
This checkpoint records the current mathematical status of the causal
information-bill programme after the sharp second-order result and the
subsequent higher-order attack.

The sharp second-order theorem is treated as established: for equal-entropy
finite sources with unequal varentropy, the exact causal residual bill is
\[
E_n
=
\sqrt{\frac{2}{\pi}}\,
|\sigma_P-\sigma_Q|\,\sqrt n
+
o(\sqrt n).
\]
The present note then develops an explicit higher-order test family, the
dyadic parity ladder.  If the first unmatched information moment is of
order \(r\), an entirely atomic B-spline buffer construction yields the
constructive causal upper bound
\[
E_n=O(n^{1/r}).
\]
For this family the exponent arises from an exact finite-difference
identity rather than from an Edgeworth heuristic or an artificial block
startup.

A matching \(n^{1/r}\) lower bound is rigorous for stationary independent
information reservoirs, but this architecture is too narrow.  Allowing
history-dependent output-correlated buffers produces a much subtler
zero-repair viability problem.  Exact finite LPs for the first nontrivial
third-order case \(r=3\) give
\[
\mathcal B_{11}=1.897576589\ldots<2,
\]
and several fixed mean-\(2\) root laws remain feasible to depth \(11\).
The data therefore no longer justify claiming a universal
\(\Theta(n^{1/3})\) lower law.  The decisive dichotomy is now:
either there exists an infinite history-indexed zero-repair tree with finite
root mean, or \(\mathcal B_n\to\infty\).

The final sections isolate the viability kernel, prove a compactness
criterion reducing bounded finite-horizon costs to existence of an infinite
tree, rule out finite-state stationary zero-repair implementations, and
state the precise next target.
\end{abstract}

\tableofcontents

\section{The causal bill}

Let \(P,Q\) be finite probability distributions.  Their single-symbol
surprisals are
\[
\ip_P(x):=-\log_2 P(x),
\qquad
\ip_Q(y):=-\log_2 Q(y).
\]
Write
\[
H(P)=\E\ip_P(X),
\qquad
V(P)=\Var(\ip_P(X)),
\]
and analogously for \(Q\).

We work with the history-indexed causal-tree contract C3-C.

\begin{definition}[Exact causal action tree]
A random action tree of depth \(n\) is exact for \((P,Q)\) if, for every
deterministic action word \(a_1\cdots a_n\in\{P,Q\}^n\), the labels observed
along the corresponding path have the product law prescribed by the action
word.
\end{definition}

Let
\[
C_n^{\mathrm{causal}}(P,Q)
\]
be the minimum Shannon entropy of a latent random seed generating such an
exact depth-\(n\) tree.

In the equal-entropy case
\[
H(P)=H(Q)=h,
\]
define the residual causal bill
\[
\boxed{
E_n(P,Q):=
C_n^{\mathrm{causal}}(P,Q)-nh.
}
\]

The central question is not whether exact simulation is possible; it is
how much \emph{additional} latent information is forced by exact causality.

\section{Established structural facts}

The following statements are part of the established framework.

\subsection{Prefix lower bound and subadditivity}

For all \(m,n\),
\[
C_{m+n}^{\mathrm{causal}}
\le
C_m^{\mathrm{causal}}+C_n^{\mathrm{causal}}.
\]
Hence
\[
C_\infty^{\mathrm{causal}}
:=
\lim_{n\to\infty}\frac{C_n^{\mathrm{causal}}}{n}
\]
exists by Fekete's lemma.

In the equal-entropy case,
\[
E_n\ge0,
\]
and \(E_n\) is subadditive.

\subsection{Projective closure}

There exists a single infinite, generally nonstationary exact causal tree
whose finite prefixes attain the asymptotic causal rate:
\[
\inf_T
\limsup_{n\to\infty}
\frac{H(T_{\le n})}{n}
=
C_\infty^{\mathrm{causal}}.
\]
A block-diagonal construction with increasing near-optimal finite blocks
gives an infinite projectively consistent substrate.

The stronger stationary or shift-invariant version is a separate problem.

\subsection{Stream-to-tree anti-diagonal lift}

Suppose a joint law of \(W^n,X^n\) satisfies
\[
W^n\sim Q^n,
\qquad
X^n\sim P^n,
\]
and the one-way causality condition
\[
\Law(X_{1:k}\mid W_{1:n})
=
\Law(X_{1:k}\mid W_{1:k})
\qquad
(k\le n).
\]
Then the anti-diagonal assignment of \(P\)-requests to the \(X\)-stream and
\(Q\)-requests to the reversed \(W\)-stream yields an exact causal action
tree.  Consequently
\[
C_n^{\mathrm{causal}}(P,Q)
\le
H(X^n,W^n),
\]
and in the equal-entropy case
\[
\boxed{
E_n(P,Q)
\le
H(X^n\mid W^n).
}
\]

Thus exact causal stream couplings are constructive upper certificates for
the tree problem.

\section{Sharp second-order theorem}

Assume
\[
H(P)=H(Q)=h
\]
and let
\[
\sigma_P^2=V(P),
\qquad
\sigma_Q^2=V(Q).
\]

The sharp second-order theorem is
\[
\boxed{
E_n(P,Q)
=
\sqrt{\frac{2}{\pi}}\,
|\sigma_P-\sigma_Q|\,\sqrt n
+
o(\sqrt n).
}
\]

Two mechanisms meet exactly.

\subsection{Adaptive lower bound}

For a deterministic adaptive policy \(\pi\), the transcript is a deterministic
function of the shared seed.  Atomwise,
\[
-\log_2 \PP(S=s)
\ge
I_\pi(Y^\pi(s)).
\]
Hence
\[
H(S)
\ge
\E\max_\pi I_\pi(Y^\pi).
\]

Threshold policies asymptotically generate an oscillating Brownian motion.
The optimal split-normal envelope has mean
\[
\sqrt{\frac{2}{\pi}}\,
|\sigma_P-\sigma_Q|,
\]
giving
\[
\liminf_{n\to\infty}
\frac{E_n}{\sqrt n}
\ge
\sqrt{\frac{2}{\pi}}\,
|\sigma_P-\sigma_Q|.
\]

\subsection{Sorted-arithmetic upper bound}

The upper construction sorts length-\(L\) source atoms by surprisal, reuses
a uniform arithmetic phase, and produces exact monotone coupling of
successive block information spectra.

After controlling the rare boundary defect and finite-bit exactization,
the resulting block increment \(Z_L\) satisfies
\[
\frac{\E Z_L^2}{L}
\longrightarrow
(\sigma_P-\sigma_Q)^2.
\]
A triangular-array functional CLT and uniform integrability of the running
maximum then give
\[
\limsup_{n\to\infty}
\frac{E_n}{\sqrt n}
\le
\sqrt{\frac{2}{\pi}}\,
|\sigma_P-\sigma_Q|.
\]

\section{The crash pair}

A convenient equal-entropy unequal-varentropy test pair is
\[
P=\left(\frac12,\frac14,\frac14\right),
\]
and
\[
Q=(q,q,1-2q),
\]
where \(q\) is the unique root giving
\[
H(Q)=\frac32.
\]
Numerically,
\[
q=0.41003242818198557608\ldots,
\]
\[
1-2q
=
0.17993514363602884784\ldots.
\]

Then
\[
H(P)=H(Q)=\frac32,
\]
\[
\sigma_P=\frac12,
\qquad
\sigma_Q
=
0.456450775263963\ldots,
\]
and therefore
\[
\boxed{
E_n
=
0.034747254051817864\ldots\sqrt n
+
o(\sqrt n).
}
\]

This pair is the main second-order stress test.

\section{Why the higher-order problem is different}

Suppose now
\[
H(P)=H(Q),
\qquad
V(P)=V(Q).
\]
The sharp second-order coefficient vanishes:
\[
\sqrt{\frac{2}{\pi}}|\sigma_P-\sigma_Q|=0.
\]

A first Edgeworth heuristic suggests that if the first unmatched centered
information cumulant is order \(r\), then normalized block spectra differ
at scale
\[
L^{-(r-2)/2},
\]
so the unnormalized monotone block mismatch is expected at scale
\[
L^{(3-r)/2}.
\]
A delayed single-scale block construction would then balance
\[
L
\quad\text{against}\quad
\sqrt n\,L^{(2-r)/2},
\]
producing
\[
L\asymp n^{1/r}
\]
and therefore the candidate residual scale
\[
n^{1/r}.
\]

This heuristic is useful but \emph{not} the basis of the rigorous
higher-order upper bound below.  In particular, lattice-phase corrections
can dominate naive Cornish--Fisher constants.

\section{The dyadic parity ladder}

Fix \(R\ge1\), and put
\[
r:=R+1.
\]

Define the surprisal laws
\[
\mu_{P_R}
=
2^{-R}
\sum_{\substack{0\le j\le R+1\\ j\ {\rm even}}}
\binom{R+1}{j}\delta_{R+j},
\]
and
\[
\mu_{Q_R}
=
2^{-R}
\sum_{\substack{0\le j\le R+1\\ j\ {\rm odd}}}
\binom{R+1}{j}\delta_{R+j}.
\]

They are realized by finite atomic probability vectors because a surprisal
level \(R+j\) corresponds to atom probability \(2^{-(R+j)}\), and the
multiplicities are integers.

Their generating functions satisfy
\[
G_{P_R}(z)-G_{Q_R}(z)
=
2^{-R}z^R(1-z)^{R+1}.
\]

Equivalently, up to translation,
\[
\boxed{
\mu_{P_R}-\mu_{Q_R}
=
2^{-(r-1)}\Delta^r.
}
\]

\begin{proposition}[Moment matching]
The first \(r-1\) information moments agree, and the first unmatched moment
is order \(r\).
\end{proposition}

\begin{proof}
The difference of generating functions contains the factor
\((1-z)^r\), so the first \(r-1\) derivatives agree at \(z=1\), while the
\(r\)th derivative does not.
\end{proof}

\subsection{The first nontrivial higher-order case}

For \(R=2\), hence \(r=3\),
\[
\boxed{
\mu_P
=
\frac14\delta_2+\frac34\delta_4,
\qquad
\mu_Q
=
\frac34\delta_3+\frac14\delta_5.
}
\]

Their common information mean is
\[
h=\frac72,
\]
and their common varentropy is
\[
V=\frac34.
\]

The centered third cumulants are
\[
\cumul_3(P)=-\frac34,
\qquad
\cumul_3(Q)=+\frac34.
\]

Thus this is an exact
\[
H_P=H_Q,\qquad V_P=V_Q,\qquad \cumul_3(P)\ne\cumul_3(Q)
\]
crash pair.

\section{Lattice phase is real}

For \(R=2\), the centered block information laws are explicit:
\[
A_L
=
-\frac{3L}{2}+2K,
\qquad
K\sim\mathrm{Bin}(L,3/4),
\]
and
\[
B_L
=
-\frac{L}{2}+2J,
\qquad
J\sim\mathrm{Bin}(L,1/4).
\]

Exact monotone transport shows two distinct subsequential \(W_2\)-limits
for even and odd \(L\).  Numerically,
\[
W_2(A_L,B_L)
\to
0.816\ldots
\]
along even \(L\), while
\[
W_2(A_L,B_L)
\to
1.054\ldots
\]
along odd \(L\).

Hence a universal third-order constant cannot be read directly from a
nonlattice Cornish--Fisher formula.  The exponent may survive, but the
constant is arithmetic.

\section{Atomic B-spline buffer}

Let
\[
U_1,\ldots,U_r
\stackrel{\rm iid}{\sim}
\mathrm{Unif}\{0,\ldots,m-1\},
\qquad
T:=U_1+\cdots+U_r.
\]
Write
\[
\PP(T=t)=\frac{c_t}{m^r}.
\]

At macrostate \(t\), create
\[
c_t2^t
\]
microatoms, each of probability
\[
\frac{1}{m^r2^t}.
\]

Thus every buffer microatom at level \(t\) has surprisal
\[
\boxed{
r\log_2m+t.
}
\]
The macro-coordinate \(t\) is therefore a true information coordinate.

\begin{proposition}[Exact buffer entropy]
The buffer entropy is
\[
\boxed{
H(B_m)
=
r\log_2m+\frac r2(m-1).
}
\]
\end{proposition}

\begin{proof}
The microatom surprisal is \(r\log_2m+T\), and
\[
\E T
=
r\,\frac{m-1}{2}.
\]
\end{proof}

The box generating function is
\[
U_m(z)
=
\frac{1-z^m}{m(1-z)}.
\]
Therefore
\[
(G_{P_R}-G_{Q_R})U_m(z)^r
=
2^{-(r-1)}m^{-r}
z^{r-1}(1-z^m)^r.
\]

\begin{theorem}[Exact atomic residual mass]
For the extended atomic laws
\[
P_R\times B_m
\qquad\text{and}\qquad
B_m\times Q_R,
\]
the common equal-surprisal atoms can be paired deterministically, and the
unmatched mass on each side is exactly
\[
\boxed{
\eps_m=m^{-r}.
}
\]
\end{theorem}

\begin{proof}
The signed residual coefficients are
\[
2^{-(r-1)}m^{-r}
(-1)^k\binom{r}{k}.
\]
The total variation of the signed residual is
\[
2m^{-r},
\]
so the positive and negative residual masses are each \(m^{-r}\).
At a fixed surprisal level, all atomic probabilities are identical;
hence the common part admits an exact atom-by-atom bijection.
\end{proof}

Conditioned on the residual output,
\[
\boxed{
H_{\rm res}
=
(r-1)+\frac r2m.
}
\]

\section{Constructive causal upper hierarchy}

Let
\[
B_0\sim B_m
\]
independently of the \(Q_R\)-stream.

At each time \(t\), use the exact atomic coupling
\[
(B_{t-1},W_t)
\longmapsto
(X_t,B_t),
\]
deterministically on the matched mass and with fresh randomness only on the
residual event.

The transition sends
\[
B_m\times Q_R
\]
exactly to
\[
P_R\times B_m.
\]
Hence
\[
(X_1,\ldots,X_n,B_n)
\sim
P_R^n\times B_m,
\]
and the construction is one-way causal.

By the chain rule,
\[
H(X^n\mid W^n)
\le
H(B_m)+n\,\eps_m H_{\rm res}.
\]

\begin{theorem}[Explicit parity-ladder upper bound]
For every integer \(m\ge1\),
\[
\boxed{
E_n(P_R,Q_R)
\le
r\log_2m
+
\frac r2(m-1)
+
\frac n{m^r}
\left(
r-1+\frac r2m
\right).
}
\]
Therefore
\[
\boxed{
E_n(P_R,Q_R)=O(n^{1/r}).
}
\]
Moreover,
\[
\boxed{
\limsup_{n\to\infty}
\frac{E_n(P_R,Q_R)}{n^{1/r}}
\le
\frac{r^2}{2}(r-1)^{-(r-1)/r}.
}
\]
\end{theorem}

\begin{proof}
The displayed finite-\(n\) bound follows from the exact initial buffer
entropy and the exact residual transition entropy.  The leading terms are
\[
\frac r2m+\frac r2\frac{n}{m^{r-1}}.
\]
Choosing
\[
m\sim((r-1)n)^{1/r}
\]
gives the stated asymptotic constant.
\end{proof}

For \(r=3\),
\[
\boxed{
\limsup_{n\to\infty}
\frac{E_n}{n^{1/3}}
\le
\frac92\,2^{-2/3}
=
2.834822362\ldots.
}
\]

This upper bound is independent of the delayed-block startup mechanism.

\section{REAL kernels and the actual smoothing hierarchy}

The REAL family contains compactly supported kernels
\[
\rho_k(x)
=
\frac{k+1}{2}(1-|x|)^k_+.
\]

For \(r=2\), the optimal discrete buffer profile is a box, the discrete
analogue of \(\rho_0\).

For \(r=3\), the exact optimizer of the reduced \(W_1\)-buffer problem is
the discrete triangular law
\[
\pi_M(s)
=
\frac{M+1-|s|}{(M+1)^2},
\qquad |s|\le M,
\]
whose continuum profile is exactly
\[
\rho_1(x)=(1-|x|)_+.
\]

Thus
\[
\boxed{
\text{REAL }\rho_1
=
\text{the exact third-order information smoother.}
}
\]

However, for \(r\ge4\), the family \((1-|x|)^k_+\) is not the correct
higher-order hierarchy.  The cusp at the origin destroys one order of
finite-difference smoothing.  The correct natural family is given by
cardinal B-splines, that is, repeated convolutions of a box.

Hence the REAL connection is exact at the first two levels but does not,
without modification, provide the full \(r\)-hierarchy.

\section{Stationary independent-buffer lower bound}

Suppose a stationary architecture has an independent information reservoir
\(B\), with surprisal
\[
J_B=-\log_2\PP(B)
\]
and
\[
H(B)=\E J_B.
\]

Then
\[
|\varphi_B(t)|
=
|\E e^{itJ_B}|
\ge
1-|t|H(B).
\]

For the parity ladder,
\[
\boxed{
|\varphi_{P_R}(t)-\varphi_{Q_R}(t)|
=
2|\sin(t/2)|^r.
}
\]

For the extended input and output information profiles
\[
A=B\times Q_R,
\qquad
Y=P_R\times B,
\]
we have
\[
|\varphi_Y(t)-\varphi_A(t)|
=
|\varphi_B(t)|
\,
|\varphi_{P_R}(t)-\varphi_{Q_R}(t)|.
\]

Kantorovich--Rubinstein yields
\[
W_1(\mu_A,\mu_Y)
\ge
\frac{
|\varphi_Y(t)-\varphi_A(t)|
}{|t|}.
\]
Choosing
\[
t\asymp \frac1{H(B)}
\]
gives
\[
W_1(\mu_A,\mu_Y)
\gtrsim
H(B)^{-(r-1)}.
\]

Using the equal-entropy information-profile converse
\[
H(Y\mid A)
\ge
\frac12W_1(\mu_A,\mu_Y),
\]
we obtain
\[
\text{cost}
\ge
H(B)
+
c_r
\frac n{H(B)^{r-1}}.
\]

Therefore, within the stationary independent-buffer architecture,
\[
\boxed{
\text{cost}=\Omega(n^{1/r}).
}
\]

Combined with the B-spline upper,
\[
\boxed{
\text{stationary independent-buffer class: }
\Theta(n^{1/r}).
}
\]

This is a genuine theorem, but it is not yet a theorem for the full C3-C
class.

\section{Why the stationary lower is too narrow}

For \(R=2\), allowing the next buffer to correlate with the current output
produces a formal stationary profile
\[
\boxed{
B_\infty(z)
=
\frac{z^2(z+1)}{3-z}.
}
\]

Its coefficients are
\[
\beta_2=\frac13,
\qquad
\beta_j=\frac4{3^{j-1}},
\quad j\ge3.
\]

At the one-step information-profile level it achieves zero repair with
finite mean.

Thus the universal statement
\[
\text{``every causal buffer must itself have size }n^{1/3}\text{''}
\]
is false.

The problem is that the child buffer law depends on the output branch, and
therefore the one-step stationary profile cannot simply be iterated while
preserving the complete iid output process.

\section{The zero-repair history-indexed tree}

For each output history
\[
h\in\bits^{\le n},
\]
let \(B_h\) be a probability law on \(\Z_{\ge0}\).

For \(R=2\), define the exact zero-repair recursion
\[
\boxed{
\mu_Q*B_h
=
\frac14\delta_2*B_{h0}
+
\frac34\delta_4*B_{h1}.
}
\tag{Z}
\]

\begin{definition}[Finite-horizon zero-repair bill]
Define
\[
\boxed{
\cB_n
:=
\inf
\left\{
\E_{B_\varnothing}J:
\{B_h\}_{|h|\le n}
\text{satisfies \eqref{Z}}
\right\}.
}
\]
\end{definition}

This object isolates the profile-level cost of exact history-compatible
information conservation.

\section{Exact finite LP data}

The recursion \eqref{Z} is linear in the buffer laws, so \(\cB_n\) can be
computed exactly as a finite LP after imposing a sufficiently large support
cutoff.

The first values are:

\[
\begin{array}{c|c|c}
n & \cB_n & \cB_n/n^{1/3}\\
\hline
1 & 0.666666667 & 0.6667\\
2 & 1.000000000 & 0.7937\\
3 & 1.185185185 & 0.8218\\
4 & 1.296296296 & 0.8165\\
5 & 1.432098765 & 0.8375\\
6 & 1.506172840 & 0.8289\\
7 & 1.596707819 & 0.8347\\
8 & 1.695473251 & 0.8477\\
9 & 1.772290809 & 0.8520\\
10& 1.846364883 & 0.8570\\
11& 1.897576589 & 0.8532
\end{array}
\]

The exact-looking rational values include
\[
\cB_{10}=\frac{1346}{729},
\qquad
\cB_{11}=\frac{4150}{2187}.
\]

The root optimum is stable when the artificial support cutoff is enlarged,
so these values are not a simple truncation artifact.

\section{The warning hidden in the same data}

The same data admit a very different interpretation.

The final gaps to \(2\) are
\[
\begin{array}{c|c}
n & 2-\cB_n\\
\hline
8&0.304526749\\
9&0.227709191\\
10&0.153635117\\
11&0.102423411
\end{array}
\]
and
\[
\frac{2-\cB_{11}}{2-\cB_{10}}
=
\frac23.
\]

Therefore the data do \emph{not} justify claiming that
\[
\cB_n\asymp n^{1/3}.
\]

The two live hypotheses are now:

\[
\boxed{
H_1:
\quad
\cB_n\to\cB_\infty<\infty,
}
\]
and
\[
\boxed{
H_2:
\quad
\cB_n\to\infty.
}
\]

The numerical range \(n\le11\) cannot decide between them.

\section{Simple mean-\(2\) roots survive deeply}

Several explicit root laws of mean exactly \(2\) remain feasible for the
zero-repair LP through depth at least \(11\).

Examples include
\[
B_\varnothing(z)
=
\left(\frac{1+z}{2}\right)^4,
\]
i.e.
\[
B_\varnothing\sim \mathrm{Bin}(4,1/2),
\]
and
\[
B_\varnothing(z)
=
\left(\frac{1+z+z^2}{3}\right)^2.
\]

Thus the hypothesis
\[
\cB_\infty\le2
\]
is not contradicted by the current finite-depth data.

However, naive greedy continuation of these roots eventually fails.
Hence finite-depth feasibility is not yet an explicit infinite construction.

\section{Greedy allocation and why it fails}

Given a parent law \(b_j=\PP(B_h=j)\), define the total input information
profile
\[
r_u
=
\frac34 b_{u-3}
+
\frac14 b_{u-5}.
\]

A child decomposition is determined by choosing a submeasure
\[
0\le c_u\le r_u,
\qquad
\sum_u c_u=\frac14,
\]
for the low-\(P\) branch.

Then
\[
b^{(0)}_j
=
4c_{j+2},
\]
and
\[
b^{(1)}_j
=
\frac43
\left(
r_{j+4}-c_{j+4}
\right).
\]

The natural greedy rule chooses the lowest information levels until
quarter mass is collected.  This works for many levels but eventually
creates an infeasible child.

For one natural geometric profile, the first failure appears at history
\[
01101101,
\]
where the child law becomes
\[
\frac23\delta_0+\frac13\delta_2.
\]

Then the next input puts mass
\[
\frac34\cdot\frac23
=
\frac12
\]
at information level \(3\), more than the low-output branch can absorb
without requiring negative information in the high branch.

Thus
\[
\boxed{
\text{greedy low-level allocation is not an infinite solution.}
}
\]

The infinite problem requires anticipative viability, not myopic matching.

\section{No finite-state stationary zero-repair transducer}

Consider an exact hidden-state transducer
\[
J_t
=
J_{t-1}+I_Q(W_t)-I_P(X_t),
\]
with
\[
W_t\mid J_{t-1}\sim Q,
\qquad
X_t\mid J_{t-1}\sim P.
\]

Because
\[
H(P)=H(Q),
\]
we have
\[
\E[J_t\mid J_{t-1}]
=
J_{t-1}.
\]
Hence \(J_t\) is a martingale.

\begin{lemma}[No stationary integrable zero-repair state]
There is no nontrivial stationary integrable state law for such a
zero-repair transducer.
\end{lemma}

\begin{proof}
Assume
\[
J_t\stackrel d=J_{t-1}
\]
and
\[
\E J_t<\infty.
\]
Take the strictly convex linearly growing function
\[
\psi(x)=\sqrt{1+x^2}.
\]
Conditional Jensen gives
\[
\E[\psi(J_t)\mid J_{t-1}]
\ge
\psi(J_{t-1}).
\]
Stationarity forces equality after expectation, hence equality holds in
strict Jensen:
\[
J_t=J_{t-1}
\quad\text{a.s.}
\]
For the \(R=2\) pair,
\[
I_Q-I_P\in\{-1,+1,+3\},
\]
so zero increment is impossible.  Contradiction.
\end{proof}

Exhaustive finite-state LP searches with up to forty state levels also
find no stationary exact zero-repair transducer.

Thus any \(H_1\) construction must be genuinely history-dependent and
multiscale.

\section{Finite automata also fail}

A stronger finite-mode ansatz allows a finite family of buffer laws
\[
B^{(1)},\ldots,B^{(K)}
\]
and deterministic output-conditioned transitions
\[
\tau_0,\tau_1:
\{1,\ldots,K\}\to\{1,\ldots,K\}
\]
s
\[
\mu_Q*B^{(s)}
=
\frac14\delta_2*B^{(\tau_0(s))}
+
\frac34\delta_4*B^{(\tau_1(s))}.
\]

Exhaustive searches find no solution for
\[
K=1,2,3.
\]

The moment mechanism suggests a general obstruction.  On a recurrent
finite component, the buffer mean is harmonic under the output transition
kernel.  A finite recurrent Markov chain forces bounded harmonic functions
to be constant.  Repeating the argument at higher moments successively
pushes the unmatched third information moment into contradiction.

This is not yet written as a complete theorem for every finite \(K\), but it
strongly supports the conclusion that any finite-state implementation is
impossible.

\section{Harvey--Peres self-similarity}

For \(R=2\),
\[
G_P(z)^2-G_Q(z)^2
=
\frac14
\left(
G_P(z^2)-G_Q(z^2)
\right).
\]

At the two-symbol information-spectrum level,
\[
\mu_P^{*2}
=
\frac1{16}\delta_4
+
\frac6{16}\delta_6
+
\frac9{16}\delta_8,
\]
and
\[
\mu_Q^{*2}
=
\frac9{16}\delta_6
+
\frac6{16}\delta_8
+
\frac1{16}\delta_{10}.
\]

The ordered common mass is \(3/4\), and the unmatched mass is \(1/4\).
After conditioning and dividing the information coordinate by \(2\), the
residual profile is again the original parity-ladder pair.

Harvey--Peres ordered matching turns this signed self-similarity into a
positive residual decomposition.  Therefore:

\[
\boxed{
\text{static positivity is not the open problem.}
}
\]

The remaining difficulty is unilateral/prefix compatibility of the future
subtrees attached to equal-total-surprisal words such as \(01\) and \(10\).

\section{Why simple Bellman reductions fail}

Several plausible lower-bound reductions have been tested.

\subsection{Terminal or prefix stochastic dominance}

Requiring only
\[
B+I_Q^{(t)}
\succeq_{\rm st}
I_P^{(t)}
\]
for all prefixes gives a bounded cost near \(0.894\), far below the full
history-indexed zero-repair LP.

Thus the obstruction is not terminal inventory.

\subsection{Depth-only dual}

A Bellman dual depending only on depth and buffer information level is
trivial.

\subsection{Count-only history compression}

Forcing the buffer to depend only on depth and the number of high-\(P\)
outputs changes the primal optimum drastically.  The ordering of the
history matters.

\subsection{Fixed finite-memory dual}

Dual certificates depending on only the final \(\ell\) output symbols
saturate for each fixed \(\ell\).

Therefore a sharp lower certificate, if one exists, must use memory across
growing scales.

\subsection{No canonical smooth cubic dual}

Raw HiGHS dual multipliers are strongly degenerate.  They do not collapse
to a canonical smooth function of
\[
j/n^{1/3}.
\]

The optimal child allocation itself develops disconnected ``islands'' at
multiple information scales, which is much more consistent with a
hierarchical self-similar certificate than with a single cubic barrier.

\section{The viability-kernel formulation}

The history labels can be eliminated conceptually.

\begin{definition}[Finite-depth viability sets]
Let
\[
\cK_0
=
\mathcal P(\Z_{\ge0}),
\]
the set of all probability laws on nonnegative integers.

Recursively, define
\[
B\in\cK_{n+1}
\]
iff there exist
\[
B_0,B_1\in\cK_n
\]
such that
\[
\boxed{
\mu_Q*B
=
\frac14\delta_2*B_0
+
\frac34\delta_4*B_1.
}
\]
\end{definition}

Then
\[
\cK_{n+1}\subseteq\cK_n,
\]
and
\[
\boxed{
\cB_n
=
\inf_{B\in\cK_n}\E_B J.
}
\]

Define the infinite viability kernel
\[
\boxed{
\cK_\infty
:=
\bigcap_{n\ge0}\cK_n.
}
\]

The dichotomy \(H_1/H_2\) becomes:

\[
\boxed{
H_1:
\quad
\exists B\in\cK_\infty
\text{with }\E_BJ<\infty,
}
\]
versus
\[
\boxed{
H_2:
\quad
\cK_\infty
\text{contains no finite-mean law.}
}
\]

\section{Compactness lemma}

The viability formulation yields a useful existence reduction.

\begin{lemma}[Bounded finite horizons imply an infinite tree]
Suppose
\[
\sup_n \cB_n<\infty.
\]
Then there exists an infinite zero-repair history-indexed tree
\[
\{B_h\}_{h\in\bits^{<\infty}}
\]
s finite root mean satisfying \eqref{Z} at every node.
\end{lemma}

\begin{proof}
Choose depth-\(n\) feasible trees with root means uniformly bounded by some
constant \(C\).

For every fixed history \(h\), let
\[
p_h
=
\left(\frac14\right)^{N_0(h)}
\left(\frac34\right)^{N_1(h)}
>0
\]
be its \(P\)-history probability.

The mean-buffer martingale identity implies
\[
\sum_{|h|=t}p_h\,\E J_h
=
\E J_\varnothing
\le C.
\]
Hence for each fixed history
\[
\E J_h
\le
\frac C{p_h}.
\]

Therefore, for each fixed \(h\), the family of laws appearing at node \(h\)
over all sufficiently deep finite trees is tight by Markov's inequality.
Since the binary tree is countable, a diagonal subsequence gives weak limits
for every fixed node.

Convolution by the finite laws \(\mu_P,\mu_Q\) is continuous under weak
convergence, and the linear identity \eqref{Z} is closed.  Therefore the
limits satisfy \eqref{Z} at every node.  The root limit has mean at most
\(C\) by lower semicontinuity.
\end{proof}

\begin{corollary}
To prove \(H_1\), it is enough to establish a uniform finite-horizon bound
\[
\boxed{
\cB_n\le C
\quad
\text{for all }n.
}
\]
To prove \(H_2\), it is enough to show
\[
\boxed{
\cB_n\to\infty.
}
\]
\end{corollary}

This is currently the cleanest decision problem.

\section{Why the old \(n^{1/3}\) renormalization target is premature}

The self-similar two-block residual suggests a formal scale relation:
information is dilated by \(2\), residual mass is reduced by \(1/4\), and
two-symbol blocking contributes another factor \(2\) in causal time.

This points to the candidate scale transformation
\[
n\mapsto8n,
\qquad
J\mapsto2J.
\]

If one could prove
\[
\cB_{8n}\gtrsim2\cB_n,
\]
then \(n^{1/3}\) would follow.

But the current finite data no longer justify assuming such a recursion.
If \(\cB_n\to2\), any multiplicative lower renormalization is false.

Therefore the correct order of attack is now:

\begin{enumerate}[label=\arabic*.]
\item decide whether \(\cK_\infty\) contains a finite-mean law;
\item only if the answer is negative, return to a quantitative
renormalization lower bound and determine the divergence exponent.
\end{enumerate}

\section{Prefix-code interpretation}

The \(R=2\) pair has a concrete binary-prefix representation.

For \(P\), take one codeword of length \(2\) and twelve codewords of length
\(4\):
\[
2^{-2}+12\,2^{-4}=1.
\]

For \(Q\), take six codewords of length \(3\) and eight codewords of length
\(5\):
\[
6\,2^{-3}+8\,2^{-5}=1.
\]

The length decisions may be coupled through the same first two fair bits:
\[
00:
\qquad
L_P=2,\quad L_Q=5,
\]
and
\[
\neq00:
\qquad
L_P=4,\quad L_Q=3.
\]

Thus the third-order parity-ladder problem can be viewed as
\[
\boxed{
\text{causal synchronization of two complete prefix parsers}
}
\]
with identical average code length
\[
\E L_P=\E L_Q=3.5,
\]
identical variance, but distinct third-order information structure.

This interpretation connects the zero-repair buffer directly to
finitary-code constructions.

\section{Relation to finitary coding}

Harvey and Peres prove that, for equal-entropy Bernoulli shifts, finite
square-root coding-length moment forces equality of informational variance.
This is the coding-theoretic analogue of the sharp second-order causal
obstruction.

Their ordered measure-preserving matching lemma also converts generating
function self-similarity into a positive residual matching.  Hence the
remaining obstruction in the present programme is not static positivity
but unilateral compatibility across future histories.

Gabor's 2024 work proves moment-optimal finitary isomorphism for
equal-entropy aperiodic Markov processes, in particular iid sources:
coding-radius moments are achievable for every
\[
\theta<\frac12
\]
in full generality, up to polylogarithmic effects.

Gabor's 2025 work extends the Harvey--Peres variance obstruction to
\(\Z^d\) and proves a Schmidt-type rigidity statement: exponentially
decaying coding-radius tails at equal entropy force equality of one-site
information laws up to permutation.  The same work obtains equality of all
information moments under exponential-tail coding.

These results do not settle the present finite-mean zero-repair viability
question, but they strongly suggest that any \(H_1\) construction must be
heavy-tailed and genuinely multiscale rather than finite-state or
exponential-tail.

\section{Relation to causal optimal transport}

Discrete-time causal optimal transport admits dynamic programming and
duality formulations.  General adapted-transport dual attainment is also
known.

Therefore the difficulty of the current lower problem is not the existence
of a dual formulation.  The difficulty is identifying an explicit
multiscale dual or viability invariant adapted to the exact parity-ladder
finite-difference self-similarity.

Recent continuous-time causal-transport work formulates analogous problems
as state-constrained stochastic control and nonlinear master equations,
reinforcing the interpretation of the buffer as an information inventory
state.

\section{Computational audit}

All computational checks used in this attack are deterministic; no Monte
Carlo evidence is used as proof.

The executable audits include:

\begin{itemize}
\item exact finite construction of sorted-arithmetic causal couplings;
\item exhaustive action-path verification for small horizons;
\item exact monotone block transport;
\item dyadic finite-bit exactization audits;
\item exact parity-ladder moment and generating-function identities;
\item global zero-repair history-tree LPs;
\item finite-support stability checks;
\item finite-state and finite-memory crash tests;
\item exact atomic B-spline residual enumeration;
\item Fourier stationary-buffer lower certificates.
\end{itemize}

The role of computation is to expose and falsify finite claims.  The
asymptotic statements are accepted only when supported by analytic proof.

\section{Current theorem ledger}

\[
\boxed{
\begin{array}{p{0.62\textwidth}|c}
\textbf{Statement} & \textbf{Status}\\
\hline
\text{Projective infinite causal substrate at rate }C_\infty^{\rm causal}
& \checkmark\\[1mm]

\text{Anti-diagonal causal stream-to-tree lift}
& \checkmark\\[1mm]

E_n=
\sqrt{2/\pi}\,|\sigma_P-\sigma_Q|\sqrt n+o(\sqrt n)
& \checkmark\\[1mm]

\text{Parity-ladder first-unmatched-moment construction}
& \checkmark\\[1mm]

\mu_{P_R}-\mu_{Q_R}=2^{-(r-1)}\Delta^r
& \checkmark\\[1mm]

\text{Fully atomic B-spline }O(n^{1/r})\text{ causal upper}
& \checkmark\\[1mm]

\text{Stationary independent-buffer }\Omega(n^{1/r})
& \checkmark\\[1mm]

\text{REAL }\rho_1\text{ equals the exact third-order smoother}
& \checkmark\\[1mm]

\text{Higher REAL }\rho_k\text{ form the universal }r\text{-hierarchy}
& \times\\[1mm]

\text{Static positive residual extraction}
& \checkmark\\[1mm]

\text{Stationary integrable zero-repair state}
& \times\\[1mm]

\text{Greedy infinite zero-repair continuation}
& \times\\[1mm]

\text{Simple finite-mode zero-repair automaton}
& \times\text{ for }K\le3\\[1mm]

\text{Zero-repair LP data through depth }11
& \checkmark\\[1mm]

\sup_n\cB_n<\infty
& \textbf{open}\\[1mm]

\cB_n\to\infty
& \textbf{open}\\[1mm]

\cB_n=\Theta(n^{1/3})
& \textbf{not established}\\[1mm]

E_n=\Theta(n^{1/r})\text{ for the full C3-C parity ladder}
& \textbf{open}
\end{array}
}
\]

\section{The present decision tree}

The research programme is now sharply bifurcated.

\subsection*{Branch \(H_1\): bounded viability}

Prove that
\[
\sup_n\cB_n<\infty.
\]
By compactness, this produces an infinite history-indexed zero-repair tree
with finite root mean.

If this happens, the universal higher-order law
\[
E_n\asymp n^{1/r}
\]
is false.  The B-spline construction remains a valid upper bound, but it is
not sharp.

The most natural current candidate is
\[
\boxed{
\cB_\infty=2.
}
\]

\subsection*{Branch \(H_2\): divergent viability}

Prove
\[
\cB_n\to\infty.
\]
Only then should one search for a quantitative renormalization or multiscale
dual certificate proving a definite divergence exponent.

If the exponent is \(1/3\), the third-order hierarchy survives.  It must,
however, arise from recursive future-compatibility rather than from block
startup, static transport, or stationary inventory.

\section{Immediate next targets}

The following order is currently the most informative.

\begin{enumerate}[label=\textbf{T\arabic*.}]
\item
\textbf{Characterize low-depth viability constraints.}
Derive analytic inequalities describing
\[
\cK_1,\cK_2,\cK_3,\ldots
\]
directly in terms of low-information tail masses of \(B\).

\item
\textbf{Search for a preserved viability family.}
Find a parameterized convex family \(\mathfrak F\) of finite-mean buffer laws
such that every
\[
B\in\mathfrak F
\]
admits children
\[
B_0,B_1\in\mathfrak F.
\]
Any such invariant family proves \(H_1\).

\item
\textbf{Prove or disprove the mean-\(2\) barrier.}
Test whether there exists
\[
B\in\cK_\infty
\]
with
\[
\E J=2.
\]

\item
\textbf{If \(H_1\) fails, construct a multiscale dual.}
The fixed-memory dual experiments show that the sharp certificate cannot be
finite-memory.  A successful dual must encode the same self-similar
hierarchy as the Harvey--Peres residual matching.

\item
\textbf{Only after the zero-repair problem is settled, transfer to the full
causal entropy bill.}
A quantitative theorem is then needed to compare positive repair entropy
with the zero-repair viability functional.
\end{enumerate}

\section{Bottom line}

The higher-order attack has now produced a rigorous constructive hierarchy
of upper bounds,
\[
\boxed{
E_n(P_R,Q_R)=O(n^{1/r}),
}
\]
together with a matching theorem in the stationary independent-buffer
architecture.

What has \emph{not} survived repeated attack is the claim that causality
alone already forces
\[
\Omega(n^{1/r})
\]
for the full history-indexed problem.

The decisive mathematical object is no longer an Edgeworth coefficient or
a block-startup balance.  It is the infinite viability kernel
\[
\boxed{
\cK_\infty
=
\bigcap_{n\ge0}\cK_n.
}
\]

For the third-order parity ladder, the next binary question is therefore
simply:

\[
\boxed{
\exists B\in\cK_\infty
\text{with }
\E J<\infty
\;?
}
\]

If yes, the conjectured universal higher-order power law collapses.

If no, the failure of finite-mean viability is itself the missing
higher-order causal invariant, and the next task is to determine its
quantitative rate.

That is the current frontier.

\begin{thebibliography}{99}

\bibitem{HarveyPeres}
N.~Harvey and Y.~Peres,
\emph{An invariant of finitary codes with finite expected square root coding length},
Ergodic Theory Dynam.\ Systems \textbf{31} (2011), 77--90.

\bibitem{DelJunco1}
A.~del Junco,
\emph{Finitary codes between one-sided Bernoulli shifts},
Ergodic Theory Dynam.\ Systems.

\bibitem{DelJunco2}
A.~del Junco,
\emph{Bernoulli shifts of the same entropy are finitarily and unilaterally isomorphic},
Ergodic Theory Dynam.\ Systems.

\bibitem{Gabor2024}
U.~Gabor,
\emph{Moment-optimal finitary isomorphism for i.i.d.\ processes of equal entropy},
arXiv:2412.15425, 2024.

\bibitem{Gabor2025}
U.~Gabor,
\emph{Moments of finitary factor maps between Bernoulli processes},
arXiv:2509.06018, 2025.

\bibitem{BackhoffEtAl2017}
J.~Backhoff-Veraguas, M.~Beiglb\"ock, Y.~Lin, and A.~Zalashko,
\emph{Causal transport in discrete time and applications},
SIAM J.\ Optim.\ \textbf{27} (2017), 2528--2562;
arXiv:1606.04062.

\bibitem{KrsekPammer}
D.~Kr\v{s}ek and G.~Pammer,
\emph{General duality and dual attainment for adapted transport},
arXiv:2401.11958.

\bibitem{BackhoffEtAl2026}
J.~Backhoff, E.~Bayraktar, I.~Ekren, and A.~Zitridis,
\emph{Analytical approach to continuous-time causal optimal transport},
arXiv:2605.19978, 2026.

\end{thebibliography}

\end{document}
